\chapter{Giới thiệu chung}
\begin{CJK}{UTF8}{min}
\label{chap:introduction}

Chương này trình bày góc nhìn tổng quan về đề tài nghiên cứu, bao gồm ba nội dung chính. Đầu tiên, phần \ref{sec:motivation} đưa ra bối cảnh, động lực và lý do lựa chọn đề tài. Tiếp theo, phần \ref{sec:prob-statement} phát biểu cụ thể bài toán cần giải quyết trong dự án. Cuối cùng, phần \ref{sec:contributions} trình bày những đóng góp chính của dự án trong việc nghiên cứu và xây dựng hệ thống Self-Evolving AI Agents hỗ trợ xử lý sự cố mạng viễn thông.

\section{Đặt vấn đề}
\label{sec:motivation}

Sự phát triển nhanh chóng của các hệ thống mạng viễn thông hiện đại đã đặt ra nhiều thách thức lớn trong quá trình vận hành, giám sát và xử lý sự cố. Các mạng thế hệ mới, đặc biệt là mạng di động 5G, có quy mô ngày càng lớn, cấu trúc ngày càng phức tạp và bao gồm nhiều thành phần có quan hệ phụ thuộc chặt chẽ với nhau như trạm gốc, thiết bị vô tuyến, lõi mạng, cấu hình truyền dẫn, tham số handover và các chỉ số hiệu năng mạng. Trong môi trường đó, sự cố có thể xuất hiện từ nhiều nguyên nhân khác nhau, từ lỗi phần cứng, sai cấu hình, nhiễu tín hiệu, suy giảm vùng phủ cho đến quá tải tài nguyên. Nếu không được phát hiện và xử lý kịp thời, các sự cố này có thể làm suy giảm chất lượng dịch vụ, gây gián đoạn kết nối, ảnh hưởng trực tiếp đến trải nghiệm người dùng và làm tăng chi phí vận hành của nhà mạng.

Trong thực tế, xử lý sự cố mạng viễn thông không chỉ dừng lại ở việc phát hiện dấu hiệu bất thường, mà còn yêu cầu xác định nguyên nhân gốc rễ đứng sau các hiện tượng quan sát được. Đây là bài toán khó vì dữ liệu vận hành mạng thường có kích thước lớn, nhiều chiều, thay đổi liên tục theo thời gian và chứa các mối quan hệ nhân quả phức tạp giữa nhiều thành phần. Các phương pháp truyền thống thường dựa vào luật thủ công, kinh nghiệm của chuyên gia hoặc các mô hình học máy tĩnh. Tuy nhiên, những phương pháp này gặp hạn chế khi hệ thống mạng thay đổi, khi xuất hiện các dạng lỗi mới hoặc khi cần đưa ra lời giải thích rõ ràng cho quá trình chẩn đoán.

Trong những năm gần đây, các mô hình ngôn ngữ lớn và AI Agent mở ra nhiều tiềm năng mới cho bài toán xử lý sự cố nhờ khả năng phân tích dữ liệu, suy luận theo ngữ cảnh, sử dụng công cụ và sinh giải thích bằng ngôn ngữ tự nhiên. Tuy vậy, các agent thông thường vẫn mang tính tĩnh, tức là hoạt động chủ yếu dựa trên tri thức, prompt, bộ nhớ hoặc chiến lược xử lý được thiết kế sẵn. Điều này khiến chúng khó thích nghi lâu dài với môi trường mạng luôn thay đổi. Từ đó, nhu cầu xây dựng các \textbf{Self-Evolving AI Agents} trở nên cần thiết. Đây là các agent có khả năng tự cải thiện dựa trên dữ liệu mới, phản hồi từ người dùng, kết quả xử lý trong quá khứ và kinh nghiệm tích lũy trong quá trình tương tác.

Chính nhu cầu đó đã trở thành động lực chính cho nghiên cứu này. Trong dự án này, chúng tôi tập trung nghiên cứu và xây dựng một hệ thống AI Agent có khả năng tự cải thiện nhằm hỗ trợ phân tích và xử lý sự cố mạng viễn thông. Hệ thống hướng tới khả năng học liên tục từ dữ liệu vận hành, cập nhật bộ nhớ, tự phản tư sau mỗi lần xử lý và sử dụng các cơ chế phản hồi để cải thiện chất lượng chẩn đoán theo thời gian. Mục tiêu cuối cùng là hỗ trợ kỹ sư mạng trong việc phát hiện sự cố, phân tích nguyên nhân gốc rễ và đề xuất hướng xử lý phù hợp, đồng thời đảm bảo kết quả đưa ra có tính giải thích và có thể hỗ trợ quá trình ra quyết định trong thực tế.

\section{Phát biểu bài toán}
\label{sec:prob-statement}

Mục tiêu của dự án là xây dựng một hệ thống \textbf{Self-Evolving AI Agents} hỗ trợ xử lý sự cố mạng viễn thông, trong đó agent có khả năng tiếp nhận dữ liệu vận hành mạng, phân tích các dấu hiệu bất thường, xác định nguyên nhân có khả năng cao nhất và đề xuất hướng xử lý phù hợp. Khác với các hệ thống chẩn đoán tĩnh, hệ thống được đề xuất cần có khả năng cải thiện theo thời gian thông qua việc học từ lịch sử xử lý, phản hồi của người dùng và kết quả đánh giá sau mỗi lần tương tác.

Cụ thể, bài toán được xem xét trong dự án này bao gồm ba nhóm nhiệm vụ chính. Thứ nhất, hệ thống cần phân tích dữ liệu đầu vào mô tả trạng thái mạng, bao gồm các chỉ số hiệu năng, thông tin cảnh báo, tham số cấu hình và các dữ liệu liên quan đến thiết bị hoặc topology mạng. Thứ hai, hệ thống cần thực hiện suy luận để xác định nguyên nhân gốc rễ của sự cố dựa trên các triệu chứng quan sát được. Thứ ba, hệ thống cần ghi nhận kết quả xử lý, phản hồi và các kinh nghiệm phát sinh để cập nhật bộ nhớ hoặc chiến lược suy luận, từ đó nâng cao hiệu quả xử lý trong các tình huống tương tự về sau.

Đầu vào và đầu ra của bài toán được định nghĩa như sau:

\textbf{Đầu vào:} Dữ liệu vận hành và giám sát mạng viễn thông, có thể bao gồm log hệ thống, cảnh báo mạng, chỉ số hiệu năng, thông tin cấu hình thiết bị, dữ liệu topology, lịch sử xử lý sự cố và phản hồi từ người dùng hoặc chuyên gia.

\textbf{Đầu ra:} Hệ thống AI Agent có khả năng hỗ trợ kỹ sư mạng thực hiện các tác vụ chính: phát hiện hoặc nhận diện sự cố, phân tích nguyên nhân gốc rễ, đưa ra giải thích cho quá trình suy luận, đề xuất hướng khắc phục và cập nhật kinh nghiệm xử lý nhằm cải thiện hiệu quả trong các lần chẩn đoán tiếp theo.

\section{Đóng góp của dự án}
\label{sec:contributions}

Đóng góp chính của dự án nằm ở việc nghiên cứu và xây dựng một hướng tiếp cận sử dụng \textbf{Self-Evolving AI Agents} cho bài toán xử lý sự cố mạng viễn thông. Thay vì chỉ xây dựng một mô hình chẩn đoán tĩnh, dự án hướng tới một hệ thống agent có khả năng học hỏi từ quá trình vận hành, phản hồi và lịch sử xử lý để từng bước cải thiện năng lực suy luận. Cụ thể hơn, các đóng góp chính của dự án được mô tả như sau:

\begin{itemize}
\item Nghiên cứu tổng quan bài toán xử lý sự cố mạng viễn thông, bao gồm các nhiệm vụ quan trọng như phát hiện sự cố, phân tích nguyên nhân gốc rễ, đánh giá tác động và đề xuất hướng khắc phục. Qua đó, dự án làm rõ những đặc điểm khiến bài toán này trở nên khó, bao gồm dữ liệu nhiều chiều, quan hệ phụ thuộc phức tạp giữa các thành phần mạng và yêu cầu cao về tính chính xác, khả năng giải thích.

\item Khảo sát và hệ thống hóa các cơ chế tự cải thiện của AI Agent, bao gồm vòng lặp phản hồi, học tăng cường, tự phản tư, cập nhật bộ nhớ và học từ lịch sử tương tác. Những cơ chế này là nền tảng để xây dựng agent có khả năng thích nghi với các tình huống mới, thay vì chỉ hoạt động dựa trên tri thức hoặc chiến lược cố định ban đầu.

\item Đề xuất kiến trúc hệ thống Self-Evolving AI Agent cho xử lý sự cố mạng viễn thông, trong đó agent có thể tiếp nhận dữ liệu vận hành, thực hiện suy luận chẩn đoán, sinh giải thích, ghi nhận phản hồi và cập nhật kinh nghiệm xử lý. Kiến trúc này hướng tới việc kết hợp khả năng suy luận của mô hình ngôn ngữ lớn với bộ nhớ, công cụ hỗ trợ và cơ chế đánh giá nhằm nâng cao chất lượng xử lý theo thời gian.

\item Xây dựng quy trình thực nghiệm và đánh giá hiệu quả của hệ thống trên các tác vụ xử lý sự cố mạng. Các tiêu chí đánh giá có thể bao gồm độ chính xác trong xác định nguyên nhân gốc rễ, chất lượng giải thích, khả năng cải thiện sau phản hồi, mức độ thích nghi với tình huống mới và hiệu quả hỗ trợ người dùng trong quá trình ra quyết định.

\item Chỉ ra các thách thức còn tồn tại khi áp dụng Self-Evolving AI Agents vào môi trường mạng viễn thông thực tế, chẳng hạn như chất lượng dữ liệu vận hành, độ tin cậy của phản hồi, rủi ro học sai từ kinh nghiệm quá khứ, khả năng kiểm soát quá trình tự cập nhật và yêu cầu đảm bảo an toàn trong các hệ thống quan trọng.

\end{itemize}

\end{CJK}
